\begin{abstract}
    
O Programa de Aquisição de Alimentos é um projeto do governo federal que realiza compras diretas de alimentos dos agricultores familiares a preço de mercado sem a necessidade de licitações e distribui para pessoas em situação de insegurança alimentar, assim como a rede socioassistencial, equipamentos públicos e a rede pública de ensino. A fim de realizar a gestão do programa no assentamento PDS - Osvaldo de Oliveira, foi desenvolvido um sistema web que contemple os processos realizados pelos agricultores e o gestor responsável. Para o desenvolvimento do sistema foram usadas abordagens de engenharia de software, como levantamento de requisitos e metodologia ágil. O projeto segue o padrão arquitetural MVC, sendo o módulo responsável pelo processamento e armazenamento de dados e o módulo de interação com o usuário separados em containers. O deploy do projeto foi feito usando o GitLab Actions para a geração da imagem e o envio para o servidor de homologação. Foram usados o Docker container para isolar os módulos e o Docker Compose para inicializar a aplicação. Todo o projeto é de código aberto e está disponível no GitLab.

\vspace{\baselineskip}
\noindent\textbf{Palavras-Chave:} Tecnologia social, Software Livre, Gestão comunitária, Engenharia de software.

\end{abstract}
